\subsection{Tipologie di utenti}

All'interno di SmartOrder esistono due tipi di utenti, ciascuno con funzionalit\`a diverse.

\subsubsection{Utente standard}

L'utente standard \`e il normale utilizzatore della piattaforma.\\
Puoi registrarti in autonomia e, una volta effettuato l'accesso, utilizzare tutte le funzionalit\`a principali per creare e gestire i tuoi ordini.

In particolare, puoi:
\begin{itemize}[itemsep=5pt, parsep=5pt, label=$\scriptstyle\bullet$]
	\item Interagire con la chat per aggiungere prodotti
	\item Visualizzare e aggiornare il carrello
	\item Inviare ordini
	\item Consultare lo storico dei tuoi ordini
	\item Accedere e gestire il tuo profilo personale
\end{itemize}

\subsubsection{Amministratore}

L'amministratore ha accesso a funzionalit\`a dedicate al controllo e alla consultazione degli ordini.

In particolare, pu\`o:
\begin{itemize}[itemsep=5pt, parsep=5pt, label=$\scriptstyle\bullet$]
	\item Visualizzare tutti gli ordini effettuati dagli utenti
	\item Filtrare gli ordini per data
\end{itemize}

A differenza dell'utente standard, l'amministratore non pu\`o registrarsi autonomamente: questo tipo di account viene creato direttamente dal sistema.

Si precisa che la registrazione \`e disponibile solo per gli utenti standard. Non \`e possibile creare un account amministratore tramite l'applicazione.

\subsection{Accesso al sistema (Autenticazione)}

\subsubsection{Pagina iniziale}

Quando apri l'applicazione, vieni automaticamente portato alla schermata iniziale, da cui puoi accedere al sistema oppure creare un nuovo account.

In questa schermata trovi due opzioni:
\begin{itemize}[itemsep=5pt, parsep=5pt, label=$\scriptstyle\bullet$]
	\item \textbf{Accedi}, se hai gi\`a un account
	\item \textbf{Registrati}, se vuoi crearne uno nuovo
\end{itemize}

Puoi passare facilmente da una modalit\`a all'altra direttamente dalla stessa pagina, senza dover cambiare schermata.

\subsubsection{Login (Accesso)}

Se hai gi\`a un account, puoi accedere selezionando la voce \vr{Accedi} nella schermata iniziale, inserendo nei campi corrispondenti le tue credenziali:
\begin{itemize}[itemsep=5pt, parsep=5pt, label=$\scriptstyle\bullet$]
	\item Username
	\item Password
\end{itemize}

\begin{figure}[H]
	\centering
	\includegraphics[width=0.8\textwidth]{ManualeUtente/login_commentato}
	\caption{Schermata di login}
\end{figure}

Dopo aver compilato i campi, premi il pulsante \vr{Accedi} per entrare nel sistema.\\
Se i dati sono corretti, verrai reindirizzato automaticamente alla schermata principale.\\
Nel campo password \`e presente l’icona (1) che ti permette di mostrare o nascondere ci\`o che stai scrivendo, utile per evitare errori durante l'inserimento.


\subsubsection{Gestione errori login}

Se inserisci credenziali non corrette, il sistema non ti permette di accedere e mostra un messaggio di errore che ti avvisa del problema.

\begin{figure}[H]
	\centering
	\includegraphics[width=0.8\textwidth]{ManualeUtente/login_errato}
	\caption{Messaggio di errore in caso di credenziali non corrette}
\end{figure}

Ad esempio:
\begin{itemize}[itemsep=5pt, parsep=5pt, label=$\scriptstyle\bullet$]
	\item Username o password errati
\end{itemize}

Inoltre, se si prova ad accedere senza compilare tutti i campi richiesti, il sistema ti segnala che i dati sono incompleti e ti invita a completarli.


Questi controlli ti aiutano a capire subito cosa non va e a correggere rapidamente l'errore.

\subsubsection{Registrazione (Creazione account)}

Se non hai ancora un account, puoi crearne uno selezionando la voce \vr{Registrati} nella schermata iniziale.

\begin{figure}[H]
	\centering
	\includegraphics[height=0.6\textheight,keepaspectratio]{ManualeUtente/registrazione_commentato}
	\caption{Schermata di registrazione}
\end{figure}

Per completare la registrazione ti verr\`a chiesto di inserire:
\begin{itemize}[itemsep=5pt, parsep=5pt, label=$\scriptstyle\bullet$]
	\item Username
	\item Email
	\item Password
	\item Conferma password
\end{itemize}

Dopo aver compilato tutti i campi, premi il pulsante \vr{Registrati} per creare il tuo account.\\
Se i dati inseriti sono corretti, la registrazione viene completata e potrai accedere al sistema.

Nei campi della password \`e presente un'icona (1) che ti permette di visualizzare ci\`o che stai digitando, cos\`i da controllare pi\`u facilmente eventuali errori.

\subsubsection{Validazione dati registrazione}

Durante la registrazione, il sistema controlla automaticamente che i dati inseriti siano corretti.\\
Se qualcosa non va, viene mostrato un messaggio che ti aiuta a capire cosa correggere.

\begin{figure}[H]
	\centering
	\includegraphics[height=0.6\textheight,keepaspectratio]{ManualeUtente/registrazione_errore_form}
	\caption{Errori dovuti alla compilazione incompleta o non valida del form}
\end{figure}

\begin{itemize}[itemsep=5pt, parsep=5pt, label=$\scriptstyle\bullet$]
	\item \textbf{Campi obbligatori}: tutti i campi devono essere compilati. Se ne lasci qualcuno vuoto, il sistema ti segnala l'errore.
\end{itemize}

\begin{itemize}[itemsep=5pt, parsep=5pt, label=$\scriptstyle\bullet$]
	\item \textbf{Requisiti della password}: la password deve rispettare alcune regole per garantire la sicurezza dell'account. Deve quindi:
	\begin{itemize}[itemsep=5pt, parsep=5pt, label=$\scriptstyle\bullet$]
		\item Avere una lunghezza tra 8 e 24 caratteri
		\item Contenere almeno una lettera maiuscola
		\item Contenere almeno una lettera minuscola
		\item Contenere almeno un numero
		\item Contenere almeno un carattere speciale
	\end{itemize}
\end{itemize}

\begin{itemize}[itemsep=5pt, parsep=5pt, label=$\scriptstyle\bullet$]
	\item \textbf{Conferma password}: il campo di conferma deve coincidere con la password inserita. Se i due valori sono diversi, il sistema ti chieder\`a di correggerli.
\end{itemize}

\begin{itemize}[itemsep=5pt, parsep=5pt, label=$\scriptstyle\bullet$]
	\item \textbf{Username esistente}: se viene inserito uno username gi\`a esistente, il sistema mostra un messaggio di errore dedicato.
\end{itemize}

Questi controlli ti guidano passo dopo passo nella compilazione, aiutandoti a creare un account valido senza errori.
